
\textbf{Problem 2.6} \rubric{Total points: 22}

Let $\cO$ denote the open sets in $\bbR^d$.
\begin{enumerate}[label=(\alph*)]
\item \rubric{ Total points: 2 \hfill

\begin{tabular}{p{0.8\linewidth}p{0.15\linewidth}}
\hline
Observe that $(a_1,b_1) \times \dots \times (a_d, b_d)$ is open & 1pt\\
Use inclusion property Lemma 2.5 & 1pt\\
\hline
\end{tabular}

}


Note that the set $(a_1,b_1) \times \dots \times (a_d, b_d)$ is open for any $a_i < b_i \in \bbR$.

 Hence $\cA_1 \subset \cA_1^\prime \subset \cO$ and thus by Lemma 2.5 we have that $\sigma(\cA_1) \subset \sigma(\cA_1^\prime) \subset \sigma(\cO) = \cB_\bbR$.
\item \rubric{ Total points: 2 \hfill

\begin{tabular}{p{0.8\linewidth}p{0.15\linewidth}}
\hline
Observe that $\supset$ is trivial & 1pt\\
Construct a rectangle $(a_1,b_1) \times \dots \times (a_d, b_d)$ inside a ball for other conclusion  & 1pt\\
\hline
\end{tabular}

}

The inclusion $\supset$ is trivial. So assume that $x \in O$. Then by definition there exist an $r > 0$ such that the ball $B_x(r) \subset O$. Now note that for any such $d$-dimensional ball, we can construct intervals $(a_i,b_i)$ such that $I := (a_1,b_1) \times \dots \times (a_d, b_d) \subset B_x(r)$ and $x \in I$. Thus $x \in \bigcup_{I \in \cA, I \subset O} I$.
\item \rubric{ Total points: 3 \hfill

\begin{tabular}{p{0.8\linewidth}p{0.15\linewidth}}
\hline
Observe that it suffices to show that open sets are in $\sigma(\cA)$ & 1pt\\
Use 2 to write $O$ as union of rectangles  & 1pt\\
Observe that $\bbQ$ is countable to conclude & 1pt\\
\hline
\end{tabular}

}

Take $O \in \cO$. If we can show that $O \in \sigma(\cA)$ then $\cB_\bbR = \sigma(\cO) \subset \sigma(\cA)$. The result then follows from 1. 

From 2 it follows that $O$ is a union over a subset collection of interval $(a_1,b_1) \times \dots \times (a_d, b_d)$ where $a_i,b_i \in \bbQ$. Since $\bbQ$ is countable and $d$ finite, the collection $\{(a_1,b_1) \times \dots \times (a_d, b_d) \, : \, a_i < b_i \in \bbQ\}$ is also countable and hence $O = \bigcup_{I \in \cA, I \subset O} I \in \sigma(\cA)$, from which it follows that $\cB_\bbR \subset \sigma(\cA)$.
\item \rubric{ Total points: 1 \hfill

\begin{tabular}{p{0.8\linewidth}p{0.15\linewidth}}
\hline
Use 1 and 3 to conclude & 1pt\\
\hline
\end{tabular}

}

This follows immediately from 1 and 3 since these imply that $\cB_\bbR = \sigma(\cA_1) \subset \sigma(\cA_1^\prime) \subset \cB_\bbR$.
\item \rubric{ Total points: 3 \hfill

\begin{tabular}{p{0.8\linewidth}p{0.15\linewidth}}
\hline
Conclude $\subset$ inclusion & 1pt\\
Assume the contrary and observe that $x_i > b_i$ must hold for some $i$ & 1pt\\
Construct an appropriate $j$ such that $x_i \notin (a_i,b_i+1/j)$ and conclude  & 1pt\\
\hline
\end{tabular}

} 

The inclusion $\subset$ is trivial, since $(a,b] \subset (a + b +1/j)$ for any $a < b \in \bbR$ and $j \in \bbN$. For the other inclusion we argue by contradiction. Suppose that 
\[
	x \in \bigcap_{j \in \bbN} (a_1, b_1+\frac{1}{j}) \times \dots \times (a_d, b_d+\frac{1}{j})
\] 
but $x \notin (a_1,b_1] \times \dots \times (a_d, b_d]$. Then, writing $x = (x_1, \dots, x_d)$, it must hold that $x_i > b_i$ for some $1 \le i \le d$ and hence there exists a $j \in \bbN$ such that $(b_i-x_i) > 1/j$. But this implies that $x_i \notin (a_i,b_i+1/j)$ which is a contradiction. So we conclude that 
\[
	(a_1,b_1] \times \dots \times (a_d, b_d] \supset \bigcap_{j \in \bbN} (a_1, b_1+\frac{1}{j}) \times \dots \times (a_d, b_d+\frac{1}{j}).
\]
\item \rubric{ Total points: 3 \hfill

\begin{tabular}{p{0.8\linewidth}p{0.15\linewidth}}
\hline
Conclude $\supset$ inclusion & 1pt\\
Use openness to find $j$ such that $b_i - (x_i+1/j) > 0$ for all $i$  & 1pt\\
Conclude & 1pt\\
\hline
\end{tabular}

}  

This time the inclusion $\supset$ is trivial since $(a,b-1/j] \subset (a,b)$ holds for every $a < b \in \bbR$ and $j \in \bbN$. For the other inclusion suppose that $x \in (a_1,b_1) \times \dots \times (a_d, b_d)$. Then, writing again $x = (x_1, \dots, x_d)$, there exists a $j \in \bbN$ such that 
\[
	x + 1/j := (x_1 + 1/j, \dots. x_d + 1/j) \in (a_1,b_1) \times \dots \times (a_d, b_d).
\] 
In particular, this implies that $b_i - (x_i+1/j) > 0$ for every $1 \le i \le d$, which implies that $x \in (a_1,b_1 - 1/j] \times \dots \times (a_d, b_d-1/j]$ and hence 
\[
	x \in \bigcup_{j \in \bbN} (a_1, b_1-\frac{1}{j}] \times \dots \times (a_d, b_d-\frac{1}{j}].
\]
\item \rubric{ Total points: 4 \hfill

\begin{tabular}{p{0.8\linewidth}p{0.15\linewidth}}
\hline
Use (e) to write every interval as countable intersection of $(a_1,b_1+1/j)$  & 1pt\\
Use (d), $\cA_2 \subset \cA_2^\prime$ and Lemma 2.2 to conclude & 1 pt\\
Use (f) to write interval as union $(a_1,b_1-1/j]$ and conclude interval is in $\sigma(\cA_2)$  & 1pt\\
Use (c) to conclude & 1pt\\
\hline
\end{tabular}

} 

It is clear that $\cA_2 \subset \cA_2^\prime$. By 5 it follows that any interval $(a_1,b_1] \times \dots \times (a_d, b_d]$ can be obtained as a countable intersection of intervals of the form $(a_1,b_1+1/j) \times \dots \times (a_d, b_d+1/j)$. By 4 $\cB_\bbR = \sigma(\cA_1^\prime)$ which by Lemma 2.1.2 contains 
\[
	\bigcap_{j \in \bbN} (a_1,b_1+1/j) \times \dots \times (a_d, b_d+1/j) = (a_1,b_1] \times \dots \times (a_d, b_d]. 
\]
So we conclude that any interval $(a_1,b_1] \times \dots \times (a_d, b_d] \in \cB_\bbR$ from which it now follows that
\[
	\sigma(\cA_2) \subset \sigma(\cA_2^\prime) \subset \sigma(\cA_1^\prime) = \cB_\bbR.
\]

For the other inclusion we consider a set $I := (a_1,b_1) \times \dots \times (a_d, b_d)$ with $a_i < b_i \in \bbQ$. Then by 6 we have that $I = \bigcup_{j \in \bbN} (a_1,b_1-1/j] \times \dots \times (a_d, b_d-1/j]$ where the later is a countable union of sets $(a_1,c_1] \times \dots \times [a_d,c_d]$ with $a_i < c_i \in \bbQ$ which must be in $\sigma(\cA_2)$ by definition of a \sigalg/. Hence, any such interval $I \in \sigma(\cA_2)$ and we thus conclude, using 3, that
\[
	\cB_\bbR = \sigma(\cA_1) \subset \sigma(\cA_2) \subset \sigma(\cA_2^\prime) \subset \sigma(\cA_1^\prime) = \cB_\bbR,
\]
which implies the result.
\item \rubric{ Total points: 2 \hfill

\begin{tabular}{p{0.8\linewidth}p{0.15\linewidth}}
\hline
Identify step to conclude $\subset$  & 1pt\\
Identify step to conclude $\supset$ & 1pt\\
\hline
\end{tabular}

} 

(2 points) Step 1 is to show that any interval $[a_1,b_1) \times \dots \times [a_d,b_d)$ can be obtained as a countable intersection of intervals $(a_1-1/j,b_1) \times \dots \times (a_d-1/j, b_d)$. From this we can conclude that any such must be in $\cB_\bbR$ proving inclusions $\subset$.

For the other inclusions we have to show that any interval $(a_1,b_1) \times \dots \times (a_d, b_d)$ can be obtained as a countable union of intervals $[a_1+1/j, b_1) \times \dots \times [a_d+1/j, b_d)$, which implies that such sets must be in the \sigalg/ generated by $[a_1,b_1) \times \dots \times [a_d, b_d)$. 
\item \rubric{ Total points: 2 \hfill

\begin{tabular}{p{0.8\linewidth}p{0.15\linewidth}}
\hline
Identify step to conclude $\subset$  & 1pt\\
Identify step to conclude $\supset$ & 1pt\\
\hline
\end{tabular}

} 
The main tool is to show that each of the intervals in these families can be obtained by taking a collection of allowed set operations for \sigalgs/, i.e. countable unions/intersections and finite complements. This will help use prove the $\subset$ inclusions. 

Then we show that any set from one of the first three families can also be obtained through countable unions/intersections and finite complements of sets from the last four families. These will then yield the $\supset$ inclusions and finish the proof.
\end{enumerate}

%\textbf{Problem 2.6} (23 points)
%Let $\cO$ denote the open sets in $\bbR$.
%\begin{enumerate}[label=(\alph*)]
%\item (2 points) Note that the interval $(a,b)$ is open for any $a < b \in \bbR$. Hence $\cA_1 \subset \cA_1^\prime \subset \cO$ and thus by Lemma 2.1.5 we have that $\sigma(\cA_1) \subset \sigma(\cA_1^\prime) \subset \sigma(\cO) = \cB_\bbR$.
%\item (2 points) The inclusion $\supset$ is trivial. So assume that $x \in O$. Then by definition there exist an $r > 0$ such that the ball $B_x(r) \subset O$. But $B_x(r) = (x-r, x+r) \in \cA_1$ so $x \in \bigcup_{I \in \cA, I \subset O} I$.
%\item (3 points) Take $O \in \cO$. If we can show that $O \in \sigma(\cA)$ then $\cB_\bbR = \sigma(\cO) \subset \sigma(\cA)$. The result then follows from 1. 
%
%From 2 it follows that $O$ is a union over a subset collection of interval $(a,b)$ where $a,b \in \bbQ$. Since $\bbQ$ is countable, the collection $\{(a,b) \, : \, a < b \in \bbQ\}$ is also countable and hence $O = \bigcup_{I \in \cA, I \subset O} I \in \sigma(\cA)$, from which it follows that $\cB_\bbR \subset \sigma(\cA)$.
%\item (1 point) This follows immediately from 1 and 3 since these imply that $\cB_\bbR = \sigma(\cA_1) \subset \sigma(\cA_1^\prime) \subset \cB_\bbR$.
%\item (3 points) The inclusion $\subset$ is trivial, since $(a,b] \subset (a + b +1/j)$ for any $j \in \bbN$. For the other inclusion we argue by contradiction. Suppose that $x \in \bigcap_{j \in \bbN} (a,b+1/j)$ but $x \notin (a,b]$. Then $x > b$ and hence there exists a $j \in \bbN$ such that $(b-x) > 1/j$. But this implies that $x \notin (a,b+1/j)$ which is a contradiction. So we conclude that $(a,b] \supset \bigcap_{j \in \bbN} (a,b+1/j)$.
%\item (3 points) This time the inclusion $\supset$ is trivial since $(a,b-1/j] \subset (a,b)$ for every $j \in \bbN$. For the other inclusion suppose that $x \in (a,b)$. Then there exists a $r > 0$ such that the interval $(x-r,x+r) \subset (a,b)$. In particular, this implies that $b - (x+r) > 0$. Now take any $j \in \bbN$ such that $j >1/(b - (x+r))$. Then $b - x > r + 1/j$ which implies that $(x-r,x+r) \subset (x-r,b-1/j]$ and hence $x \in \bigcup_{j \in \bbN} (a,b-1/j]$.
%\item (4 points) It is clear that $\cA_2 \subset \cA_2^\prime$. By 5 it follows that any interval $(a,b]$ can be obtained as a countable intersection of intervals of the form $(a,b+1/j)$. By 4 $\cB_\bbR = \sigma(\cA_1^\prime)$ which by Lemma 2.1.2 contains $\bigcap_{j \in \bbN} (a,b+1/j) = (a,b]$. So we conclude that any interval $(a,b] \in \cB_\bbR$ from which it now follows that
%\[
%	\sigma(\cA_2) \subset \sigma(\cA_2^\prime) \subset \sigma(\cA_1^\prime) = \cB_\bbR.
%\]
%
%For the other inclusion we consider a set $(a,b)$ with $a,b \in \bbQ$. Then by 6 we have that $(a,b) = \bigcup_{j \in \bbN} (a,b-1/j]$ where the later is a countable union of sets $(c,d]$ with $c,d \in \bbQ$ which must be in $\sigma(\cA_2)$ by definition of a \sigalg/. Hence, any interval $(a,b) \in \sigma(\cA_2)$ and we thus conclude, using 3, that
%\[
%	\cB_\bbR = \sigma(\cA_1) \subset \sigma(\cA_2) \subset \sigma(\cA_2^\prime) \subset \sigma(\cA_1^\prime) = \cB_\bbR,
%\]
%which implies the result.
%\item (2 points) Step 1 is to show that any interval $[a,b)$ can be obtained as a countable intersection of intervals $(a-1/j,b)$. From this we can conclude that any set $[a,b)$ must be in $\cB_\bbR$ proving inclusions $\subset$.
%
%For the other inclusions we have to show that any interval $(a,b)$ can be obtained as a countable union of intervals $[a+1/j, b)$, which implies that $(a,b)$ must be in the \sigalg/ generated by $[a,b)$. 
%\item (3 points) The main tool is to show that each of the intervals $(-\infty, a],$ $(-\infty ,a), (a,\infty)$ and $[a,\infty)$ can be obtained by taking any allowed set operation for \sigalgs/, i.e. countable unions/intersections and finite complements. This will help use prove the $\subset$ inclusions. 
%
%Then we show that any set of the form $(a,b), [a,b)$ or $(a,b]$ can also be obtained through countable unions/intersections and finite complements of intervals of the forms $(-\infty, a]$, $(-\infty ,a), (a,\infty)$ and $[a,\infty)$. These will then yield the $\supset$ inclusions and finish the proof.
%\end{enumerate}

\bigskip

\textbf{Problem 2.9}

First note that if $\mu(A \cap B) = \infty$ then by property 2 we have that also $\mu(A), \mu(B)$ and $\mu(A \cup B) = \infty$ and hence the result holds trivially. So assume now that $\mu(A \cap B) < \infty$. Since 
\[
	A \cup B = (A \setminus (A\cap B)) \cup (B \setminus (A \cap B)) \cup (A \cap B),
\] 
it follows from property 1 that
\[
	\mu(A \cup B) = \mu(A \setminus (A \cap B))) + \mu(A \cap B) + \mu(B \setminus (A \cap B)).
\] 
Adding $\mu(A \cap B) < \infty$ to both side we get
\begin{align*}
	\mu(A \cup B) + \mu(A \cap B) 
	&= \mu(A \setminus (A \cap B))) + \mu(A \cap B) + \mu(B \setminus (A \cap B)) + \mu(A \cap B)\\
	&= \mu(A) + \mu(B), 
\end{align*}
where the last line follows from applying property 3 twice.  
 
\bigskip 

\textbf{Problem 2.10} \rubric{Total points: 15}

\begin{enumerate}[label=(\alph*)]
\item \rubric{ Total points: 5 \hfill

\begin{tabular}{p{0.8\linewidth}p{0.15\linewidth}}
\hline
Define sets $E_i$ & 1pt\\
Check that they are disjoint & 1pt\\
Prove union for $A_i$ property (by induction) & 2pt\\
Use previous part to prove equality of unions & 1 pt\\
\hline
\end{tabular}

}

Define the sets $E_i$ as follows:
\[
	E_1 = A_1, \quad E_k = A_k \setminus A_{k-1} \text{ for } k \ge 2.
\]

Clearly, $E_k \cap E_t = \emptyset$ for any $k \ne t$.

For the other properties, note that $A_i = \bigcup_{k = 1}^i E_k$ holds for $i = 1$. Suppose now that it holds for a given $i$. Then
\begin{align*}
	A_{i + 1} &= A_{i} \cup (A_{i+1} \setminus A_i) \\
	&= \left(\bigcup_{k = 1}^i E_k\right) \cup (A_{i+1} \setminus A_i)\\
	&= \left(\bigcup_{k = 1}^i E_k\right) \cup E_{i+1} = \bigcup_{k = 1}^{i+1} E_k,
\end{align*}
and thus $A_i = \bigcup_{k = 1}^i E_k$ holds for all $i \ge 1$ by induction. 

This then immediately implies that $\bigcup_{i \ge 1} A_i = \bigcup_{k \ge 1} E_k$.
\item \rubric{ Total points: 3 \hfill

\begin{tabular}{p{0.8\linewidth}p{0.15\linewidth}}
\hline
Use (a) and $\sigma$-additivity to write measure of union $A_i$ as sum measure $E_k$ & 2pt\\
Use this to conclude & 1pt\\
\hline
\end{tabular}

}


By a) and $\sigma$-additivity we have that
\begin{align*}
	\mu(A_i) &= \mu\left(\bigcup_{k = 1}^i E_k\right) = \sum_{k = 1}^i \mu(E_k),
\end{align*}
and 
\[
	\mu\left(\bigcup_{i \ge 1} A_i\right) = \mu\left(\bigcup_{k \ge 1} E_k\right) = \sum_{k = 1}^\infty \mu(E_k).
\]

Thus,
\[
	\lim_{i \to \infty} \mu(A_i) = \sum_{k = 1}^\infty \mu(E_k) = \mu\left(\bigcup_{i \ge 1} A_i\right),
\]
as required.
\item \rubric{ Total points: 7 \hfill

\begin{tabular}{p{0.8\linewidth}p{0.15\linewidth}}
\hline
Define sets $B_i$& 1pt\\
Use part 1 and Proposition 2.12 to write limit of measures as measure of union & 1pt\\
Use finiteness measure $A_1$ and exclusion property to write limit measures $B_i$ as difference measures $A_1$ and limit & 2pt\\
Write measure union $B_i$ as difference measure $A_1$ and intersection $A_i$ & 2pt\\
Conclude & 1pt\\
\hline
\end{tabular}

}

Consider the family $B_i = A_1 \setminus A_i$. 

Then $(B_i)_{i \ge 1}$ is an increasing family of sets satisfying $\bigcup_{i \ge 1} B_i = A_1$ and thus, by part 1 of Proposition 2.12 we have
\[
	\lim_{i \to \infty} \mu(B_i) = \mu\left(\bigcup_{i \ge 1} B_i\right).
\]

Since $\mu(A_1) < \infty$, the exclusion property (point 3 in Proposition 2.11) gives us that
\[
	\lim_{i \to \infty} \mu(B_i) = \mu(A_1) - \lim_{i \to \infty} \mu(A_i). 
\]

On the other hand
\[
	\bigcup_{i \ge 1} B_i = \bigcup_{i \ge 1} (A_1 \setminus A_i) = A_1 \setminus \bigcap_{i \ge 1} A_i
\]
and thus
\[
	\mu\left(\bigcup_{i \ge 1} B_i\right) = \mu\left(A_1 \setminus \bigcap_{i \ge 1} A_i\right)
	= \mu(A_1) - \mu\left(\bigcap_{i \ge 1} A_i\right),
\]
where we used the exclusion property since $ \mu\left(\bigcap_{i \ge 1} A_i\right) \le \mu(A_1) < \infty$.

Together we get
\[
	\mu(A_1) - \lim_{i \to \infty} \mu(A_i) = \mu(A_1) - \mu\left(\bigcap_{i \ge 1} A_i\right),
\]
which yields the required identity since $\mu(A_1) < \infty$.

\end{enumerate}

\bigskip

\textbf{Problem 2.11}

The idea is to construct a family of disjoint sets $(E_i)_{i \in \bbN}$ with the following properties:
\begin{enumerate}
\item $E_i \subset A_i$, and
\item $\bigcup_{i \in \bbN} E_i = \bigcup_{i \in \bbN} A_i$.
\end{enumerate}

If such a sequence exists then we have
\begin{align*}
	\mu(\bigcup_{i \in \bbN} A_i) &= \mu(\bigcup_{i \in \bbN} E_i) &&\text{by 2}\\
	&= \sum_{i = 1}^\infty \mu(A_i) &&\text{because $E_i$ are disjoint and $\mu$ is $\sigma$-additive}\\
	&\le \sum_{i = 1}^\infty \mu(A_i) &&\text{by 1 and monotone property of $\mu$}.
\end{align*}

So we are left to construct the required family of sets $(E_i)_{i \in \bbN}$. The following set will do:
\[
	E_1 = A_1 \quad E_i = A_i \setminus \bigcup_{k < i}^i A_k \text{ for all } i > 1.
\]
Note that by definition the set $E_i$ are pair-wise disjoint and property1 holds. Finally, property 2 holds since $\bigcup_{i = 1}^k E_i = \bigcup_{i = 1}^k A_i$ holds for all $k \ge 1$.

\bigskip

\textbf{Problem 2.13}

\begin{enumerate}[label=(\alph*)]
\item

We first make the following observations about $\mathcal{N}$:
\begin{itemize}
\item because $\mu(\emptyset) = 0$ it holds that $\emptyset \in \mathcal{N}$,
\item if $N, M \in \mathcal{N}$ then $N \setminus M \in \mathcal{N}$ since $N \setminus M \subset N$, and
\item if $(N_i)_{i \ge 1}$ is a family of sets in $\mathcal{N}$ then so is $\bigcup_{i \ge 1} N_i$.
\end{itemize}

From the first point it follows that $\emptyset = \emptyset \cup \emptyset \in \overline{\cF}$ and $\Omega = \Omega \cup \emptyset \in \overline{\cF}$.

Furthermore, if $A, B \in \cF$ and $N, M \in \mathcal{N}$, then by the second point and because $A \setminus B \in \cF$,
\[
	(A \cup N) \setminus (B \cup M) = (A \setminus B) \cup (N \setminus M) \in \overline{\cF}.
\]

Finally, let $(A_i \cup N_i)_{i \ge 1}$ be a collection of sets in $\mathcal{N}$. Then using the third point we get
\[
	\bigcup_{i \ge 1} A_i \cup N_i = \bigcup_{i \ge 1} A_i \cup \bigcup_{i \ge 1} N_i \in \overline{\cF}.
\]
\item (1 pt) From the definition we immediately get that $\mu(\emptyset) = 0$. Now, let $(A_i \cup N_i)_{i \ge 1}$ be a collection of disjoint sets in $\mathcal{N}$. Then
\[
	\bar{\mu}(\bigcup_{i \ge 1} A_i \cup N_i) = \bar{\mu}(\bigcup_{i \ge 1} A_i \cup \bigcup_{i \ge 1} N_i)
	= \mu(\bigcup_{i \ge 1} A_i) = \sum_{i \ge 1} \mu(A_i) = \sum_{i \ge 1} \bar{\mu}(A_i \cup N_i).
\]
\item 
This follows from the fact that $\bar{\mu}|_\cF(A) = \bar{\mu}(A \cup \emptyset) = \mu(A)$.
\item 
Suppose that $N \subset \Omega$ is a null set for $\overline{\cF}$. Then there exists an $A \cup M \in \overline{\cF}$ such that $N \subset A \cup M$ and $\bar{\mu}(A \cup M) = \mu(A) = 0$. However, since $M \in \mathcal{N}$, there must also exist a $B \in \cF$ with $M \subset B$ and $\mu(B) = 0$. But this implies that $N \subset A \cup B \in \cF$ which implies that $N \in \mathcal{N}$. Therefore, since $N = \emptyset \cup N$ it follows that $N \in \overline{\cF}$ and hence every null set is part of the \sigalg/.
\end{enumerate}